\subsection*{6.5. A reminder on the Picard--Lefschetz formula (XV 3.4)}

Let $S$ be the spectrum of a strictly henselian valuation ring, with
generic point $\eta$, geometric generic point $\bar\eta$, and closed
point $s$.  Let
\[
\rho:X\longrightarrow S
\]
be a proper morphism, where $X$ is a connected regular scheme of
dimension $n$.  Suppose that $\rho$ is smooth away from a single point
$x$ in the special fiber and that $x$ is an ordinary quadratic
singularity of that fiber.

Let $I$ denote the inertia group of $S$ at $s$.  By XV 3.4, $I$ acts
trivially on $H^i(X_{\bar\eta},\mathbb Q_\ell)$ for $i\ne n-1$, and its
action on $H^{n-1}(X_{\bar\eta},\mathbb Q_\ell)$ is described as follows.

\subsubsection*{6.5.1.}
If $n=2m+1$, then $I$ acts through a quotient isomorphic to
$\mathbb Z/2\mathbb Z$ (if $\operatorname{char}(k)\ne2$, there is only
one such quotient).  There is an element
\[
\delta\in H^{n-1}(X_{\bar\eta},\mathbb Q_\ell(m)),
\]
determined uniquely up to sign and called the \underline{vanishing
cycle at $s$}, such that
\begin{equation}\tag{6.5.1.1}
\delta\cdot\delta=(-1)^m2
\quad\text{in}\quad
H^{2n-2}(X_{\bar\eta},\mathbb Q_\ell(n-1))\simeq\mathbb Q_\ell.
\end{equation}
The nontrivial element $\sigma$ of $\mathbb Z/2\mathbb Z$ acts on
$H^{n-1}(X_{\bar\eta},\mathbb Q_\ell)$ by the reflection
\begin{equation}\tag{6.5.1.2}
\sigma x=x-(-1)^m(x\cdot\delta)\delta.
\end{equation}
\footnote{The print numbers the reflection formula (6.5.2.2), although
it is the second formula of 6.5.1 and is followed by (6.5.2.1).}
Here $x\cdot\delta$ belongs to
$H^{2n-2}(X_{\bar\eta},\mathbb Q_\ell(m))\simeq\mathbb Q_\ell(-m)$,
so $(x\cdot\delta)\delta$ defines an element of
$H^{n-1}(X_{\bar\eta},\mathbb Q_\ell)$.

\subsubsection*{6.5.2.}
If $n=2m+2$, then $I$ acts through its canonical quotient
$\mathbb Z_\ell(1)$.  There is an element
\[
\delta\in H^{n-1}(X_{\bar\eta},\mathbb Q_\ell(m)),
\]
canonically determined up to sign and called the \underline{vanishing
cycle at $s$}, such that $\sigma\in\mathbb Z_\ell(1)$ acts on
$x\in H^{n-1}(X_{\bar\eta},\mathbb Q_\ell)$ by the transvection
\begin{equation}\tag{6.5.2.1}
\sigma(x)=x-(-1)^m\sigma\,(x\cdot\delta)\delta.
\end{equation}
Indeed, $x\cdot\delta$ belongs to
$H^{2n-2}(X_{\bar\eta},\mathbb Q_\ell(m))\simeq\mathbb Q_\ell(-m-1)$.
Multiplication by $\sigma\in\mathbb Z_\ell(1)$ therefore gives an
element $\sigma(x\cdot\delta)\in\mathbb Q_\ell(-m)$, whose product with
$\delta\in H^{n-1}(X_{\bar\eta},\mathbb Q_\ell(m))$ lies in
$H^{n-1}(X_{\bar\eta},\mathbb Q_\ell)$.

We return to the global situation.

\subsubsection*{Theorem 6.6.}
Let $X$ be a smooth connected projective variety of dimension $n$ over
an algebraically closed field $k$.  Suppose that $n$ is even or that
$\operatorname{char}(k)\ne2$.  Let $D$ be a Lefschetz pencil for $X$,
let $\eta$ be its generic point, let $\bar\eta$ be a geometric point
over $\eta$, and let $s_1,\ldots,s_N$ be the points of
$D\cap\check X$.

For each $i$, choose an isomorphism $\ell_i$ between $\bar\eta$ and the
geometric generic point of $S_i$, the strict henselization of $D$ at
$s_i$.  Via $\ell_i$, let $I_i$ be the inertia group at $s_i$ acting on
$H^{n-1}(X_{\bar\eta},\mathbb Q_\ell(j))$, and let $\delta_i$ be the
vanishing cycle at $s_i$ defined in 6.5.\footnote{The print calls
$I_i$ the inertia group ``of $\delta_i$'' before $\delta_i$ is defined.
The standard construction is the inertia group at $s_i$, embedded by
the chosen path $\ell_i$.}
Thus
\[
\delta_i\in H^{n-1}(X_{\bar\eta},\mathbb Q_\ell(m)),
\qquad
m=\begin{cases}
(n-1)/2,& n\text{ odd},\\
n/2-1,& n\text{ even}.
\end{cases}
\]
\footnote{The print omits the generic-fiber subscript in this display;
the preceding action and the definition in 6.5 place every
$\delta_i$ in the cohomology of $X_{\bar\eta}$.}
Choose the $\ell_i$ so that the subgroups
\[
I_i\subset
\pi=\pi_1^{\mathrm{mod}}(D-D\cap\check X,\bar\eta)
\]
topologically generate $\pi$, as is possible by SGA 1, XIII 2.10.
Then:

\medskip
\noindent
\textup{6.6.1.} If $X$ satisfies condition (LV), the vanishing
cohomology space
\[
E^{n-1}(X_{\bar\eta},\mathbb Q_\ell(m))
\stackrel{\mathrm{def}}{=}
H^{n-1}(X,\mathbb Q_\ell(m))^\perp
\subset H^{n-1}(X_{\bar\eta},\mathbb Q_\ell(m))
\]
(cf. 5.4.4) is generated by the vanishing cycles
$\delta_i$, $i=1,\ldots,N$.

\medskip
\noindent
\textup{6.6.2.} The vanishing cycles $\delta_i$ are conjugate, up to
sign, under the global monodromy group
$\pi_1(D-D\cap\check X)$.

\subsubsection*{Proof.}
For 6.6.1, the Picard--Lefschetz formula 6.5 and the fact that the
$I_i$ generate $\pi=\pi_1(D-D\cap\check X)$ give
\begin{equation}\tag{6.6.3}
H^{n-1}(X_{\bar\eta},\mathbb Q_\ell(m))^\pi
=\bigl\langle\delta_1,\ldots,\delta_N\bigr\rangle^\perp
\quad\text{in }H^{n-1}(X_{\bar\eta},\mathbb Q_\ell(m)).
\end{equation}
By 5.6.10,
\begin{equation}\tag{6.6.4}
H^{n-1}(X_{\bar\eta},\mathbb Q_\ell(m))^\pi
=\operatorname{Im}\bigl(
H^{n-1}(X,\mathbb Q_\ell(m))
\longrightarrow H^{n-1}(X_{\bar\eta},\mathbb Q_\ell(m))
\bigr).
\end{equation}
Taking orthogonals in (6.6.3) proves 6.6.1.

For 6.6.2, note simply that the vanishing cycle $\delta_i$ is
determined up to sign by the representation of the inertia group at
$s_i$.  Apply 6.1.2, which gives conjugacy in
$\pi_1^{\mathrm{mod}}(\check{\mathbb P}^{r}-\check X,\bar\eta)$ of the
composite inertia homomorphisms
\begin{center}
\begin{tikzcd}[column sep=huge,row sep=huge]
\mathbb Z_\ell(1) \arrow[r]
  & \pi_1^{\mathrm{mod}}(D-D\cap\check X,\bar\eta) \arrow[d] & {}\\
{}
  & \pi_1^{\mathrm{mod}}(\check{\mathbb P}^{r}-\check X,\bar\eta) & .
\end{tikzcd}
\end{center}
Consequently, the vanishing cycles determined by these homomorphisms
are conjugate up to sign by elements of
$\pi_1^{\mathrm{mod}}(\check{\mathbb P}^{r}-\check X,\bar\eta)$.
The proof is completed by 6.1.1, which gives the surjection
\[
\pi_1^{\mathrm{mod}}(D-D\cap\check X)
\longrightarrow
\pi_1^{\mathrm{mod}}(\check{\mathbb P}^{r}-\check X).
\]

